% Auto-generated by build_fig_workflow.py
% DO NOT hand-edit coordinates; edit the Python script instead.
\documentclass[tikz,border=0pt]{standalone}
\usepackage{ctex}
\usepackage{tikz}
\usetikzlibrary{positioning,fit,backgrounds,calc,arrows.meta}
\begin{document}
\begin{tikzpicture}[
  >=Stealth,
  node distance=0cm,
  cell/.style={{
    draw=black, line width=0.5pt, rounded corners=3pt, fill=white,
    minimum width=3.300cm, minimum height=1.400cm,
    align=center, font=\footnotesize, inner sep=2.5pt, text width=2.900cm
  }},
  title/.style={{
    font=\footnotesize\bfseries, text=black, align=left, anchor=south west
  }},
  bar/.style={{
    draw=black, line width=0.5pt, rounded corners=3pt, fill=gray!30,
    minimum width=10.700cm, minimum height=0.650cm,
    align=center, font=\footnotesize, text width=10.300cm
  }},
  smalllabel/.style={{
    font=\scriptsize, fill=white, inner sep=2pt, rounded corners=2pt,
    draw=black, line width=0.4pt
  }},
  arrow/.style={{
    ->, line width=0.7pt, color=black
  }},
  dasharrow/.style={{
    ->, line width=0.7pt, color=black, dashed
  }}
]

\node[cell, minimum width=2.400cm, text width=2.000cm] (n0_0) at (-3.500,2.700) {\textbf{选择实验}\\ Lab（1–4）};
\node[cell, minimum width=2.400cm, text width=2.000cm] (n0_1) at (-3.500,0.700) {\textbf{提交案情 / 启动研判}};
\node[cell, minimum width=2.400cm, text width=2.000cm] (n0_2) at (-3.500,-1.300) {\textbf{失败场景诚实反思}};
\node[cell, minimum width=3.300cm, text width=2.900cm] (n1_0) at (0.000,2.700) {\textbf{预处理}\\ 案情清洗 / 字段对齐};
\node[cell, minimum width=3.300cm, text width=2.900cm] (n1_1) at (0.000,0.700) {\textbf{BGE 嵌入 + HAN 构图}\\ 语义 / 结构双通道};
\node[cell, minimum width=3.300cm, text width=2.900cm] (n1_2) at (0.000,-1.300) {\textbf{LangGraph 反思闭环}\\[-1pt]\scriptsize plan→preprocess\\[-3pt]\scriptsize analyze→cluster→reflect};
\node[cell, minimum width=2.800cm, text width=2.400cm] (n2_0) at (3.500,2.700) {\textbf{规则 vs AI 建议对比}\\ Lab1 / Lab2};
\node[cell, minimum width=2.800cm, text width=2.400cm] (n2_1) at (3.500,0.700) {\textbf{冻卡决策与伦理权衡}\\ Lab3};
\node[cell, minimum width=2.800cm, text width=2.400cm] (n2_2) at (3.500,-1.300) {\textbf{边界认知与诚实反思}\\ Lab4};

\begin{scope}[on background layer]
\node[draw=black, line width=0.5pt, rounded corners=4pt, fill=gray!18, inner sep=0.400cm, fit=(n0_0) (n0_1) (n0_2)] (panel0) {};
\node[draw=black, line width=0.7pt, rounded corners=4pt, fill=gray!25, inner sep=0.400cm, fit=(n1_0) (n1_1) (n1_2)] (panel1) {};
\node[draw=black, line width=0.5pt, rounded corners=4pt, fill=gray!18, inner sep=0.400cm, fit=(n2_0) (n2_1) (n2_2)] (panel2) {};
\end{scope}
\node[title] at ([yshift=2pt]panel0.north west) {\textbf{阶段 1}\\[-4pt] 学情 / 案情输入};
\node[title] at ([yshift=2pt]panel1.north west) {\textbf{阶段 2}\\[-4pt] FraudLens 系统内核};
\node[title] at ([yshift=2pt]panel2.north west) {\textbf{阶段 3}\\[-4pt] 人机协同决策};

\draw[arrow] (n0_0.south) -- (n0_1.north);
\draw[arrow] (n0_1.south) -- (n0_2.north);
\draw[arrow] (n1_0.south) -- (n1_1.north);
\draw[arrow] (n1_1.south) -- (n1_2.north);
\draw[arrow] (n2_0.south) -- (n2_1.north);
\draw[arrow] (n2_1.south) -- (n2_2.north);
\draw[arrow] (n0_2.east) -- (n1_2.west) node[midway, above=18pt, smalllabel] {案情输入};
\draw[arrow] (n1_2.east) -- (n2_2.west) node[midway, above=18pt, smalllabel] {候选结果};
\node[bar, minimum width=10.700cm, text width=10.300cm] (audit) at (0.100,-3.475) {审计日志 → 加权评价量表 + 思政 / 伦理映射 → 回流“研—训—评”};
\draw[arrow] (n2_2.south) -- (n2_2.south |- audit.north) node[pos=0.35, left=6pt, smalllabel] {决策记录};
\coordinate (turnL) at (-5.950,-3.475);
\coordinate (turnUp) at (-5.950,-1.300);
\draw[dasharrow] (audit.west) -- (turnL) -- (turnUp) -- (n0_2.west);
\node[smalllabel, anchor=west] at (-5.850,-2.388) {评价/反思回流};

\end{tikzpicture}
\end{document}
